% Lösungen Zone (zusammenbau v0.1)
\einheitenkopf[][Kennst du schon]{Das kennst du schon}

Zuordnung: 1 → der Differenzenquotient ist genau diese Steigung, Einheit 1 bis 3; die Sekantengleichung in Einheit 1, das Steigungsdreieck an der Tangente in Einheit 2 · 2 → Zähler des Differenzenquotienten in Einheit 1, Kontrolle in Einheit 2 und 4 · 3 → jede Sachzeile in Einheit 1, 2 und 4 · 4 → f' für den Ableitungswert in Einheit 2, für f'(x) = c in Einheit 2 und 4, f'' in Einheit 4 · 5 → die Ableitung gleich einem vorgegebenen Wert setzen in Einheit 2 und 4, die zweite Ableitung null setzen in Einheit 4 · 6 → für das Anlegen der Tangente in Einheit 2 und 3 und für die Beurteilung von Aussagen am Graphen · 7 → das Werkzeug, mit dem in Einheit 4 die Rate als Funktion untersucht wird (Maximum von f' über die Nullstellen von f'' und das Vorzeichen von f''' oder den Vorzeichenwechsel; Randwerte) · 8 → der Differenzenquotient ist genau diese Steigung, Einheit 1 bis 3; die Sekantengleichung in Einheit 1, das Steigungsdreieck an der Tangente in Einheit 2 · 9 → der Differenzenquotient ist genau diese Steigung, Einheit 1 bis 3; die Sekantengleichung in Einheit 1, das Steigungsdreieck an der Tangente in Einheit 2

\erg{1}{a) $m = \frac{8 - 2}{3 - 1} = 3$ \quad b) $m = 2$ (1 nach rechts, 2 nach oben), $g\colon y = 2x - 1$}
\erg{2}{a) $f(4) = 16 - 12 = 4$ \quad b) $h(3) = 15 \cdot \mathrm{e}^{-0{,}6} \approx 8{,}23$}
\erg{3}{a) $3{,}5$ h \quad b) $2$ h und $45$ min; $\frac{9{,}2 - 8}{8} = 0{,}15 = 15\,\%$}
\erg{4}{a) $f'(x) = 3x^2 - 4$ \quad b) $h'(t) = 1 \cdot \mathrm{e}^{-0{,}5t} + t \cdot (-0{,}5) \cdot \mathrm{e}^{-0{,}5t} = (1 - 0{,}5t) \cdot \mathrm{e}^{-0{,}5t}$}
\erg{5}{a) $t = 4$ \quad b) $3t \cdot (t - 6) = 0$, also $t_1 = 0$ und $t_2 = 6$}
\erg{6}{a) rechts von $x = 1$ – dort geht der Graph nach rechts oben \quad b) Der Baum wächst anfangs immer schneller (der Graph wird steiler), etwa ab der Mitte immer langsamer (der Graph wird flacher).}
\erg{7}{a) $f'(x) = 2x - 6 = 0$, also $x = 3$ \quad b) $f'(x) = 3x^2 - 12 = 0$, also $x = 2$ oder $x = -2$; $f''(x) = 6x$, $f''(-2) = -12 < 0$: Maximum bei $x = -2$}
\erg{8}{a) Fehler: Ben nimmt nur den Höhenunterschied; die Steigung ist Höhenunterschied geteilt durch waagerechten Unterschied. Richtig: $m = \frac{11 - 3}{6 - 2} = 2$}
\erg{9}{a) $m = \frac{-1 - 7}{5 - 1} = -2$}
